\documentclass[a4paper, 10pt]{article}
\usepackage[T1, T2A]{fontenc}
\usepackage[utf8]{inputenc} 
\usepackage[english, russian]{babel}
\usepackage{amsfonts,amssymb,amsthm,mathtools,amsmath}
\usepackage{indentfirst} 
\usepackage{graphicx}
\usepackage{hyperref}
\usepackage{biblatex}
\usepackage{geometry}
\addbibresource{lit.bib}


\geometry{
	a4paper,
	total={170mm,257mm},
	left=20mm,
	top=20mm,
}
\theoremstyle{plain}
\newtheorem*{pro}{Доказательство}
\newtheorem*{ex}{Пример}
\newtheorem{lem}{Лемма}
\newtheorem{theorem}{Теорема}
\newtheorem{defin}{Определение}
\newtheorem*{rem}{Замечание}
\newtheorem*{sled}{Следствие}
\newcommand{\defeq}{:=}
\newcommand{\tang}[1][\hspace{0pt}]{\frac{\partial }{\partial #1}}
\newcommand\Lie{Lie}
\author{я}
\title{,,,}
\begin{document}
	\maketitle
		\begin{equation}
		e_1^{\circ(N)}
		=
		\sum_{k=0}^{N-1}\alpha_k(s)e_k \quad \text{где } \quad \alpha_k(0)=0.
		\label{eq:fundamental_relation} 
	\end{equation}
		\begin{equation}
		\left.\frac{\partial\alpha_k}{\partial s_r}\right|_0=0 \qquad\text{если }k\ne r.
		\label{eq:weight_diagonal}
	\end{equation}
		\begin{equation}
		e_a\circ e_b=e_{a+b}
		\qquad\text{для }a+b\le N-1.
		\label{eq:low_product}
	\end{equation}
	\begin{lem}[$A_N$]
		\[B_r=rB_1,\text{для } r=1,\ldots,N-2.
		\text{ Эквивалентно, } B_1:B_2:\cdots:B_{N-1}=1:2:\cdots:(N-1).\]
	\end{lem}
	\begin{proof}
		Пусть $m_s$ тензор умножения в точке $s=(s_0,\ldots,s_{N-1})$ и
		\[D_r:=\left.\frac{\partial m_s}{\partial s_r}\right|_{s=0}=\bigl(\Lie_{e_r}(\circ)\bigr)_0.\]
		$D_r$ симмертческое билинейное отображение \(D_r:T_0M\times T_0M\longrightarrow T_0M.\)
		Выберем $p,q\ge 1$ т.ч.\(p+q=N\), из ассоциативности и \eqref{eq:fundamental_relation},
		\[e_p\circ e_q=e_1^{\circ(N)}
		=
		\sum_{k=0}^{N-1}\alpha_k(s)e_k.\]
		Дифференцируя по $s_r$ в нуле и используя  \eqref{eq:weight_diagonal}, получаем
		\begin{equation}
			D_r(e_p,e_q)=B_r e_r,
			\label{eq:Dr_boundary}
			\quad \text{где }\lambda_r:=
			\left.\frac{\partial\beta}{\partial s_r}\right|_0.
		\end{equation}
		%Теперь, для нечетного \(N\), все \(\lambda_r\) из градуировки,  для четного же \(N\), только для одного \(r\) возможно $\lambda_r = \left.\frac{\partial\beta}{\partial s_r}\right|_0.\ne0$ \par
		Пусть теперь $a,b\ge1$ и $a+b\le N-1/$ Из \eqref{eq:low_product},
		\(e_a\circ e_b=e_{a+b}.\)
		Применим условие интегрируемости
		\[\Lie_{X\circ Y}(\circ)
		=
		X\circ\Lie_Y(\circ)+Y\circ\Lie_X(\circ)\]
		с $X=e_a$ и $Y=e_b$, и ограничим в ноль. Получаем
		\begin{equation}
			\begin{aligned}
				D_{a+b}(U,V) = \bigl(\Lie_{e_{a+b}}(U\circ V)\bigr)_0
				=&
				\bigl(e_a \circ \Lie_{e_b}(U\circ V)\bigr)_0 + \bigl(\Lie_{e_a}(U\circ V) \circ e_b\bigr)_0
				=\\=&
				e_a\circ_0D_b(U,V)
				+e_b\circ_0D_a(U,V)
				\label{eq:HM_linearized} 
				\qquad \forall U,V\in T_0M.
			\end{aligned}
		\end{equation}
		
		Возьем $U=e_p$ и $V=e_q$. По \eqref{eq:Dr_boundary}, слева получаем
		\[
		D_{a+b}(e_p,e_q)
		=B_{a+b}e_{a+b}.
		\]
		Справа,
		\begin{align*}
			e_a\circ_0D_b(e_p,e_q)
			&=e_a\circ_0(B_b e_b) = B_b e_{a+b} \\
			e_b\circ_0D_a(e_p,e_q)
			&=e_b\circ_0(B_a e_a) = B_a e_{a+b}.
		\end{align*}
		%Заметим, что 
		%\(e_c\circ_0 h=0
		%\qquad\text{for every }c\ge1,\)
		%другими словами $x^cy=0$, что следует из $2xy=0$.
		%Значит, произведения, содержащие \(\lambda\) в \eqref{eq:HM_linearized} исчезают. 
		\[
		B_{a+b} e_{a+b} = B_be_{a+b}+B_ae_{a+b} \qquad (a,b \ge 1, a+b \le N-2)
		\]
		Сравнивая коэффициенты и делая индукцию получаем
		\[
		B_r=rB_1,\qquad r=1,\ldots,N-2.
		\]
	\end{proof}
\end{document}